\subsection{5.3 Implications for reverse stress testing and further questions}\label{conclusions}\label{v29-concluding-discussion}\label{stress-test-implications}
An exact initial sensitivity need not predict how much dependence stress is required to reverse a diversification trade. The distinction is relevant to reverse sensitivity analysis \cite{pesenti2019,pesenti2022}, where the quantity of interest is the distance to a decision reversal. Under the traded-coordinate dominance condition \eqref{eq:legacy-3.4}, \(r_{0,\varepsilon}=m_\varepsilon/K_\varepsilon\) extrapolates an attained initial slope, yet if \(\gamma_1=c_{\min}<c_3\),
\[
\frac{\rho^\pm_\varepsilon}{r_{0,\varepsilon}}\longrightarrow\frac{s^*}{\gamma_1J}>1.
\]
Proposition~\ref{prop:balanced-linear-profile} attains the initial slope by balancing opposite movements on selected fibers. Equality in the efficiency bound at reversal would require balancing cheapest-coordinate projections across the entire tail section, which the corner obstruction of Theorem~\ref{thm:zero-premium} excludes under the stated cost condition. Unlike optimal allocation across full Fr\'echet or Wasserstein ambiguity sets \cite{liu2026}, our reverse stress test fixes a proposed trade and measures the least perturbation of a common reference law that reverses its paired ES ranking.

For a whole-book comparison across reference dependence structures, the size of the nominal advantage must also enter the assessment. In the family of Section~5.2, consider the boundary perturbation \(\boldsymbol\gamma(\eta)=(\gamma_0,\gamma_0,\gamma_0(1+\eta))\). The gain-normalized critical budget and its right derivative at equal effective costs satisfy
\[
\mathcal R_\vartheta(\eta):=\lim_{\varepsilon\downarrow0}
\frac{\rho^\pm_{\varepsilon,\vartheta}(\eta)}{m_{\varepsilon,\vartheta}}
=\frac{p\,s^*_\vartheta(\eta)}{J_\vartheta},\qquad
\mathcal R_\vartheta'(0+)=p\gamma_0a_{+,\vartheta}.
\]
Theorem~\ref{thm:projected-flow} gives the absolute right derivative \((s^*_\vartheta)'(0+)=\gamma_0J_\vartheta a_{+,\vartheta}\). By \eqref{eq:legacy-7.2}, increasing \(\vartheta\) raises the gain-normalized response \(\mathcal R_\vartheta'(0+)\) but lowers the absolute response: \mbox{\(\frac{d}{d\vartheta}(\gamma_0J_\vartheta a_{+,\vartheta})=\gamma_0(J_\vartheta'a_{+,\vartheta}+J_\vartheta a_{+,\vartheta}')<0\)}. Thus the decline in the nominal-gain scale outweighs the increase in background share. Reporting both boundary sensitivities to a risk committee distinguishes a greater relative need for background displacement from a smaller absolute response in budget units.

Beyond the guaranteed initial linear interval, characterizing the boundaries between the active and stationary regions of an optimal kernel could help determine the full profile \(F\) between the analytical envelopes. A uniform rate for \(\sup_{s\le S_0}|\varepsilon^{-2}D_\varepsilon(s\varepsilon^2)-F(s)|\), combined with the lower growth estimate in Theorem~\ref{thm:strict-crossing}, would yield finite-trade error bounds for the critical budgets.

In the compact-box benchmark, a restoration premium can arise from dependence geometry and full marginal conservation even with bounded marginal tails. For unbounded marginals such as integrable Pareto or lognormal laws, the inverse-CDF derivative is unbounded near probability one, so the uniform displacement estimate in the recovery argument would need a tail-weighted substitute. One-dimensional \(W_1\) remains the \(L^1\) integral of the CDF discrepancy, so this loss of a uniform quantile bound does not by itself imply a different restoration-cost exponent. The open problem is to control that integral jointly with the paired-risk error and determine whether the \(\varepsilon^2\) critical scale persists under suitable tail assumptions. Higher-dimensional projected-flow representations and estimation of critical budgets from loss data raise related questions, alongside recent work on ES estimation under weak tail assumptions \cite{bartleckstein2026}.
